\section{Mean-field analysis}

\subsection{Fixed points}
Solve $\dot{c}=\dot{s}=\dot{i}=0$ for weak/strong Allee. Classify extinction, coexistence, and saddle points. Record coordinates $(c^*,s^*,i^*)$ and feasibility ($>0$). In current scans (weak/strong, $\mu\in\{0,1\}$) all equilibria with $i^*>0$ are saddles (Re\,$\lambda_{\max}>0$); they serve as unstable seeds for spatial simulations (low-$c^*$ near 0 and high-$c^*>0.9$). Typical weak-Allee saddles: $(c^*,s^*,i^*)\approx(1.1,0.5,0.5)$ with Re\,$\lambda_{\max}\sim 60$; strong-Allee saddles cluster near $(1.0,0.0,0.1)$ with Re\,$\lambda_{\max}\sim 9$.

\subsection{Linear stability}
Compute the Jacobian $J(c^*,s^*,i^*)$; eigenvalues $\lambda_j$ determine stability. For spatial perturbations, use $\lambda_j(k)=\mathrm{eig}(J-Dk^2)$ with $D=\mathrm{diag}(D_c,D_s,D_i)$.

\subsection{Allee-induced bifurcations}
Vary $a$ and $r_c$ to track saddle-node or transcritical bifurcations separating extinction vs proliferation branches. Weak Allee yields soft thresholds; strong Allee produces sharp barriers and hysteresis-like behavior. Nonreciprocity tilts the diagrams and shifts thresholds.

\subsection{Role of nonreciprocity}
Nonreciprocal couplings shift nullclines and can tilt bifurcation diagrams. Increasing $\beta,\eta$ amplifies immune suppression and raises Re\,$\lambda_{\max}$; increasing $\delta$ and $r_i$ lowers Re\,$\lambda_{\max}$ but did not produce stability in the explored grid. Report the two families of saddles (subcritical vs over-threshold) under weak/strong Allee, $\mu=0/1$.

\subsection{Representative findings}
\begin{itemize}
  \item Weak Allee, $\mu=1$: all $i^*>0$ equilibria unstable; Re\,$\lambda_{\max}$ median $\sim 40$; $c^*$ mediana $\sim1.01$, $s^*\sim0$, $i^*\sim0.37$.
  \item Strong Allee, $\mu=1$: Re\,$\lambda_{\max}\sim9$; $s^*$ casi nulo; barrera crítica más marcada.
  \item Disminuir $r_c$ o $\beta,\eta$ y aumentar $r_d,\delta$ son direcciones para buscar estabilidad en barridos futuros.
\end{itemize}

% Fig. 2: phase plane / nullclines (weak/strong, mu=0/1). Fig. 3: bifurcation diagrams in $a$ or $r_c$.
\begin{figure}[!htbp]
    \centering
    \includegraphics[draft=false,width=0.48\linewidth]{steady_mu0.png}
    \includegraphics[draft=false,width=0.48\linewidth]{steady.png}
    \caption{Phase planes y nullclines: weak Allee con $\mu=0$ (izq.) y $\mu=1$ (der.). Los estados estacionarios son umbrales inestables.}
\end{figure}

\begin{figure}[!htbp]
    \centering
    \includegraphics[draft=false,width=0.48\linewidth]{steady_strong_mu_0.png}
    \includegraphics[draft=false,width=0.48\linewidth]{steady_strong_mu_1.png}
    \caption{Phase planes para efecto Allee fuerte con $\mu=0$ (izq.) y $\mu=1$ (der.). La barrera crítica es más abrupta y $s^*\approx 0$.}
\end{figure}

\begin{figure}[!htbp]
    \centering
    \includegraphics[draft=false,width=0.48\linewidth]{estabilidad_lineal_weak_mu0.png}
    \includegraphics[draft=false,width=0.48\linewidth]{estabilidad_lineal_weak_mu1.png}
    \caption{Espectros lineales (weak Allee): $\mu=1$ atenúa modos de alta frecuencia pero no estabiliza $k\approx 0$.}
\end{figure}

\begin{figure}[!htbp]
    \centering
    \includegraphics[draft=false,width=0.48\linewidth]{estabilidad_lineal_strong_mu0.png}
    \includegraphics[draft=false,width=0.48\linewidth]{estabilidad_lineal_strong_mu1.png}
    \caption{Espectros lineales (strong Allee): $\mu=1$ reduce la banda inestable en altas frecuencias, sin eliminar la inestabilidad en $k\approx 0$.}
\end{figure}

\clearpage

